
We began by observing that a model's alignment strategy --- DPO or SFT --- might be
inferrable from standard benchmark scores alone. We confirmed this intuition: 83.3\%
classification accuracy (permutation p=0.031) confirms that the alignment fingerprint
is present and statistically significant at pilot scale. But we also found that the
fingerprint tells a different story than expected. It is not fairness that marks
DPO-aligned models as distinct; it is truthfulness.

\subsubsection{Summary}

We introduced alignment fingerprinting as a binary classification problem over 4D
trustworthiness benchmark vectors and showed that:

\begin{enumerate}
\item **DPO and SFT 7B models are separable from public benchmark scores at 83.3\% LOO-CV
\end{enumerate}
   accuracy (permutation p=0.031).** Alignment auditing without training-data access is
   practically feasible.

\begin{enumerate}
\item **The dominant discriminative dimension is TruthfulQA MC2 (Fisher's criterion=0.8122),
\end{enumerate}
   not fairness benchmarks.** DPO models score +4.6pp higher on truthfulness on average
   across matched pairs (per-pair mean delta), directly contradicting the predicted
   fairness-advantage mechanism. BBQ shows no systematic DPO advantage (k=3/6, p=0.66).

\begin{enumerate}
\item \textbf{Misclassification cases reveal alignment space structure.} An RLHF model clustering
\end{enumerate}
   with DPO suggests preference-based methods share a benchmark signature; a DPO model
   misclassified as SFT reveals that base architecture proximity can dominate alignment
   signal for near-identical pairs.

\subsubsection{Future Directions}

\textbf{Mechanism confirmation.} Compare DPO models trained on datasets with explicit factual
quality ratings (UltraFeedback) vs. purely pairwise preference data (Anthropic HH-RLHF)
to test whether the truthfulness signal is annotation-driven.

\textbf{Architecture-controlled replication.} Restrict analysis to single-base-model families
to determine whether 83.3\% fingerprint accuracy partially reflects base model variation.

\textbf{Three-class alignment fingerprinting.} Extend to DPO vs RLHF-PPO vs pure-SFT
classification to formalize the RLHF/DPO boundary observation.

\textbf{Full-task evaluation.} Remove the 100-sample limit for publication-quality results.

\subsubsection{Closing}

Alignment fingerprinting from benchmark scores is feasible --- but it fingerprints the
unexpected. The signal that identifies DPO-trained models is not the fairness advantage
that DPO's design suggests, but a truthfulness advantage that its design does not
explicitly optimize for. This inversion should prompt revisiting what preference-based
training actually learns, and how alignment auditing tools should be designed when
theory and empirical signal diverge.

